\chapter{Conclusion}
\label{chap:conclusion}

\section{Contributions and main findings}

This thesis derives an Expanded Welch--Satterthwaite test for comparing
categorical MI across independent populations. Its mathematical
contribution is the observation-level sensitivity of the estimated MI
variance, accounting for changes in cell probabilities, margins and
pointwise MI. The variance of this sensitivity gives a first-order
sampling-variance approximation, which is used with moment matching to
choose MI-specific degrees of freedom for a Student \(t\) distribution.

The experiment evaluates the mean and sampling variance approximations
separately before comparing the completed rejection rules. It contains
366 evaluation settings, with independent moment fitting at 186 null
settings and 100,000 repetitions per batch. The first research question
is answered by the different errors in the two moment approximations:
the sampling-variance prediction can be accurate while the mean estimate
remains inflated, substantially lowering the population first-order
degrees of freedom relative to the moment target. Near independence,
the quadratic sampling contribution explains the scale of the observed
component estimates and their departure from the first-order values.

For the second research question, Expanded Welch improves false-positive
rates relative to Wald and the sample-size comparator in selected
small-sample populations, with a nominal power cost. At matched null
rejection rates, it has higher power along the selected uniform
\(4\times4\), \(n=20\), and different-skew \(8\times8\), \(n=100\),
paths. Its calibrated differences from Wald are practically small in
the different-skew \(4\times4\), \(n=10{,}000\), comparison. These
results identify a benefit of sample-dependent MI degrees of freedom
beyond an overall reduction in rejection.

The method does not resolve errors in the shared statistic. In sparse
different-skew tables, MI bias can shift or reverse the estimated
difference and cause rejection to fall initially as the true difference
increases. Same-skew populations can instead remain conservative when
the expected difference bias cancels. The matched construction and
allocation checks show why population structure and sample allocation
matter alongside the cutoff. Expanded Welch therefore offers useful
behaviour for some categorical comparisons, while accurate estimation
of the MI difference and its uncertainty remains necessary for reliable
inference.

\section{Limitations and future work}

The findings concern fixed categorical populations, the tested association
constructions and the specified alternative paths. They support further
investigation of the MI-specific rule rather than its use as a general
replacement for Normal Wald. The independently fitted references and
calibration thresholds assess mechanisms under known populations; they
do not supply an immediately available procedure for arbitrary observed
tables.

A first extension is to derive the second-order mean contribution to the
variance estimator and evaluate a corresponding degrees-of-freedom
correction. A second is to improve sparse MI bias and standard-error
estimation, assessing their effects separately through null rejection,
power curves and estimated direction. The supporting jackknife check
provides a starting point for the bias problem.

Near independence, the quadratic sampling approximation may support a
more accurate comparison of MI differences. Any likelihood-ratio or
resampling procedure must preserve the equal-MI null even when the
complete population tables differ. Reversed power paths, additional
association constructions and broader allocation comparisons would
establish how far the observed calibrated-power advantage extends.
